\section{One arbitrary head}\label{sec:onehead}

Theorem~F needs every head to have bounded owner-load. One head of arbitrary owner-load is harmless on its own: the walk
DAG of a single walk is planar, and $\TCP(1)$ holds by separators~\cite[Section~5.1]{FourthPaper2026}.
\Cref{lem:W} gives $\TCP(1)$ in the form of a label that is anchored to cells rather than to times, and Theorem~M
(\cref{thm:M}) combines it with the family clock: one arbitrary walk plus walks of bounded owner-load whose families use
boundedly many scales are depth-reducible. \Cref{sec:one-counter} collects what Theorems F and~M, and a stack
formulation of TCP, say about where a depth-robust family must live.

\subsection{Lemma W: a strip label for one walk}\label{sec:one-W}

Fix integers $w\ge2$ and $r\ge1$, a \emph{line residue} $\rho\in\{0,\dots,w-1\}$ and an \emph{offset}
$\omega\in\{0,\dots,r-1\}$. The \emph{lines} are the cells congruent to~$\rho$ modulo~$w$, and the \emph{strips} are the
sets of $w-1$ consecutive cells strictly between consecutive lines. For a walk~$x$ of length~$T$, the \emph{line times} are
the times~$t$ at which $x(t)$ is a line. For a time~$t$ that is not a line time, let $S(t)$ be the strip of~$x(t)$ and
$c(t)$ the number of times $t'\le t$ with $x(t')\in S(t)$; the \emph{strip label} is
$\Psi(t)=(c(t)+\omega)\bmod r$. The \emph{strip cut} consists of the edges of the walk DAG with an endpoint at a line time,
the hop edges $u\to p$ with $\Psi(u)\le\Psi(p)$, and the state edges $t\to t-1$ with $\Psi(t)<\Psi(t-1)$. (The value of
the label at line times does not matter, since every edge at a line time is cut.)

\begin{lemma}[Lemma W: the strip cut for one walk]\label{lem:W}
\formal{OneWalk.lean}{strip_cut, strip_cut_depth}
Let $x$ be a walk of length~$T$ and let $w\ge2$ and $r\ge1$ be integers. For every line residue and offset, every path of
the walk DAG of~$x$ that avoids the strip cut contains at most $r-1$ hop edges. For some line residue and offset, the
strip cut has at most $4T/w+wT/r$ edges.
\end{lemma}

\begin{outline}
In Lean a walk is a map $\N\to\Z$ (\lean{IsWalk}); a walk of length~$T$ is extended by staying at $x(T-1)$, and since
hops, line times, strip counts and cut edges are taken at times below~$T$, the extension changes none of them. Hops are
the pairs of consecutive visits of a cell at distance at least two in time (\lean{IsHop}, \lean{hops});
lines, strips, the strip count and the strip label are \lean{IsLine}, \lean{strip}, \lean{stripCount} and
\lean{stripLabel}, and the cut consists of \lean{cutStates} and \lean{cutHops}. \emph{Depth.} Along a path that avoids the
cut, the label does not increase along state edges and decreases strictly along hop edges, so \cref{lem:labels} bounds the
number of hop edges by $r-1$. \emph{Size.} Some residue has at most $T/w$ line times, by averaging over the $w$ residues
(\lean{exists_line_residue}). A \emph{wrap} is a time $\tau$, not a line time, with $c(\tau)+\omega\equiv0\pmod r$; for
some offset there are at most $T/r$ wraps, again by averaging (\lean{exists_offset}). A hop at a cell that is not a line has
both endpoints in one strip~$S$, and $c$ rises by at least one from $p$ to~$u$, at the visit~$u$; so if the hop is cut, the
count of~$S$ passes a multiple of~$r$ at a visit to~$S$ in $(p,u]$, a wrap of~$S$. At each wrap at most $w-1$ hops of
the strip span it, one per cell, by~(P1) (\lean{hop_unique}). Two adjacent cells that are not lines lie in one strip
(\lean{strip_eq_of_adj}), so a state edge between non-line times is cut only at a wrap. Hence at most $2L+W$ state edges
and at most $L+(w-1)W$ hop edges are cut, where $L$ and $W$ are the numbers of line times and of wraps
(\lean{card_cutStates_le}, \lean{card_cutHops_le}), and $3L+wW\le3T/w+wT/r$.
\end{outline}

For $0<\eps\le4$, with $w=\lceil8/\eps\rceil\ge2$ and $r=\lceil2w/\eps\rceil$, the strip cut has at most $\eps T$ edges
and depth at most $r-1<2w/\eps\le16/\eps^2+2/\eps$. The label is local to the strip of the current cell: the two endpoints of a hop read
the counter of the same strip. So the label is not a band rule in the sense of the fourth paper's nesting
barrier~\cite[Definition~5.5]{FourthPaper2026}, and that barrier does not apply to it.

\subsection{Theorem M}\label{sec:one-M}

Let $x_0$ be any walk of length~$T$; its long hops are called \emph{chords}. Let $H'$ be a set of long hops of further
walks $x_1,\dots,x_k$ of length~$T$, of owner-load at most~$\ell$, with families, bands and events as in
\cref{sec:nest-families}, and put $N=2k\ell$, an upper bound for the number of virtual heads (\cref{lem:families}).
Let $G$ be the DAG on $[0,T)$ with the state edges, all chords and the hop edges of~$H'$; an edge is determined by its
endpoints.

\begin{theorem}[Theorem M: one arbitrary head]\label{thm:M}
\notformalized
Let $q\ge1$ and $h\ge0$ be integers and $0\le\beta_1<\beta_2<\dots<\beta_q$ integers, and suppose that every family~$F$
of~$H'$ is assigned a window $i(F)\in\{1,\dots,q\}$ with $\beta_{i(F)}\le\lambda(d_F)\le\beta_{i(F)}+h$. Let
$\eps\in(0,1]$ and put
\[
  w=2^{\lceil\log_2(8(4+2k)/\eps)\rceil},\qquad K_0=\lceil8(w+1)/\eps\rceil,\qquad
  K=\Bigl\lceil\frac8\eps\Bigl(q+qN+\frac{64q^2N^22^h}\eps+qwN+N^22^{h+1}\Bigr)\Bigr\rceil.
\]
Then $G$ has a cut of at most $\eps T$ edges and depth less than $\DM(k,\ell,q,h,\eps)=K_0K^{2q}$. The bound does not
depend on~$T$, on $x_0$ or on the $\beta_i$, and pauses are allowed.
\end{theorem}

For $q=0$, that is, $H'=\emptyset$, \cref{lem:W} applies. Windows may touch or overlap; any cover of a finite set of bands
by $q$ windows of width~$h$ can moreover be replaced by the greedy cover, in which each window starts at the smallest band
not yet covered, with at most $q$ windows of the same width.

\begin{informalproof}
\emph{Widths and lines.} Put $W_0=w$ and, for $1\le i\le q$, let $W_i$ be the least power of~$2$ that is at least
$\max(\mathrm{lo}_i,w)$, where $\mathrm{lo}_i=64qN2^{\beta_i+h}/\eps$. The $\mathrm{lo}_i$ increase with~$i$, so
$w=W_0\mid W_1\mid\dots\mid W_q$, and $W_i\le2\max(\mathrm{lo}_i,w)\le2(\mathrm{lo}_i+w)$. Draw $\rho$ uniformly from
$\{0,\dots,W_q-1\}$. The \emph{level-$i$ lines} are the cells congruent to~$\rho$ modulo~$W_i$; every level-$i$ line is a
level-$i'$ line for $i'\le i$. A \emph{level-$i$ strip} is the set of cells strictly between two consecutive level-$i$
lines, so it has $W_i-1$ cells and every level-$i$ strip lies inside a level-$(i+1)$ strip. Let
$Z=\{t:x_0(t)\text{ is a level-}0\text{ line}\}$, and for $t\notin Z$ let $\sigma_i(t)$ be the level-$i$ strip of $x_0(t)$.

\emph{Counters and label.} The level-$0$ counter of a level-$0$ strip~$S$ is $c^0_S(t)=\#\{t'\le t:x_0(t')\in S\}$, and for
$i\ge1$ the level-$i$ counter of a level-$i$ strip~$S$ counts \emph{entries} into~$S$,
$c^i_S(t)=\#\{t'\le t:x_0(t')\in S\text{ and }x_0(t'-1)\text{ is a level-}i\text{ line}\}$. For the walks of~$H'$, let
$n_i(t)$ be the number of events at times at most~$t$ of scale at least~$\beta_i$. Choose independent uniform offsets
$\omega_0\in[0,K_0)$ and $\omega_1,\dots,\omega_q,\omega'_1,\dots,\omega'_q\in[0,K)$, and put
$d_0(t)=(c^0_{\sigma_0(t)}(t)+\omega_0)\bmod K_0$, $d_i(t)=(c^i_{\sigma_i(t)}(t)+\omega_i)\bmod K$ and
$e_i(t)=(n_i(t)+\omega'_i)\bmod K$ for $1\le i\le q$. The label $\Psi(t)$, for $t\notin Z$, is the mixed-radix number with
the \emph{static} digit order, most significant first,
\[
  d_q,\ e_q,\ d_{q-1},\ e_{q-1},\ \dots,\ d_1,\ e_1,\ d_0,
\]
so it has $K_0K^{2q}$ values; put $\Psi(t)=0$ for $t\in Z$. The cut~$F$ consists of all edges with an endpoint in~$Z$,
the hop edges $u\to p$ (chords or hops of~$H'$) with $\Psi(u)\le\Psi(p)$, and the state edges at descents. Every edge
at~$Z$ is cut, so the value of $\Psi$ on~$Z$ is never compared, and by \cref{lem:labels} $F$ has depth less than
$K_0K^{2q}$. A digit \emph{wraps} at a step if its count passes a value congruent to minus its offset modulo its radix;
a digit whose count rises by~$\Delta$ at a step wraps there with probability at most $\Delta$ divided by its radix.

\emph{The deciding level of a chord.} Let $\gamma=(p,u)$ be a chord at a cell~$c$ that is not a level-$0$ line (chords at
lines have both endpoints in~$Z$). Its \emph{excursion} $E(\gamma)=\{x_0(t):p<t<u\}$ is a nonempty interval of cells next
to~$c$, on one side of~$c$. Let $\lev(\gamma)$ be the largest $i\in\{0,\dots,q\}$ such that $E(\gamma)$ contains a level-$i$
line, and $0$ if it contains none. If $\lev(\gamma)=i<q$, then $E(\gamma)$ contains no level-$(i+1)$ line, so
$E(\gamma)\cup\{c\}\subseteq\sigma_{i+1}(c)$ (\emph{confinement}).

\emph{(M-a) State edges.} Let $t,t-1\notin Z$. Then $x_0(t-1)$ and $x_0(t)$ are equal or adjacent cells that are not lines,
so they lie in the same strip at every level. No counter of level $i\ge1$ changes at~$t$, since an entry at~$t$ needs
$x_0(t-1)$ on a line; the level-$0$ counter rises by one, and $n_i$ by the number of events at~$t$ of scale at least
$\beta_i$. So $\Psi(t)<\Psi(t-1)$ only if some digit wraps at~$t$.

\emph{(M-b) Chords.} Let $\gamma=(p,u)$ be a chord at a cell~$c$ that is not a line, with $i=\lev(\gamma)$. (i) At $p$ and
at~$u$ every head-$0$ digit reads the counter of the same strip $\sigma_{i'}(c)$, and counters do not decrease. (ii) The
level-$i$ counter of $\sigma_i(c)$ rises strictly in $(p,u]$: for $i=0$ by the visit at~$u$; for $i\ge1$, $E(\gamma)$
contains the nearest level-$i$ line~$y_*$ to~$c$ on the side of the excursion, and if $t_1$ is the last time in
$(p,u)$ with $x_0(t_1)=y_*$, then $x_0(t_1+1)$ is the neighbour of $y_*$ on the side of~$c$, since the walk reaches~$c$
at~$u$ without returning to~$y_*$; this neighbour is a cell of $\sigma_i(c)$, so $t_1+1\in(p,u]$ is an entry into
$\sigma_i(c)$. (iii) If $i<q$, then by confinement $x_0$ visits no
level-$(i+1)$ line at the times $p,\dots,u-1$, so no counter of a level greater than~$i$ changes in $(p,u]$. (iv) The event
counts do not decrease. Hence, if no digit at a position at or above~$d_i$ wraps in $(p,u]$, the most significant digit
whose count changed has risen without wrapping, and $\Psi(u)>\Psi(p)$. By~(iii), $\gamma$ is violated only if
\textup{(c1)} the level-$i$ counter of $\sigma_i(c)$ wraps at a time in $(p,u]$, or \textup{(c2)} some $e_{i'}$ with
$i'>i$ wraps in $(p,u]$.

\emph{(M-c) Hops of~$H'$.} Let $\eta=(p,u)\in H'$ lie in a family~$F$ with window $i=i(F)$, and $p,u\notin Z$. If $x_0$
visits no level-$i$ line during $(p,u)$, it stays in one level-$i$ strip on $[p,u]$, hence in one level-$i'$ strip for every
$i'\ge i$, and enters none of them in $(p,u]$; so the digits $d_{i'}$, $i'\ge i$, do not change. The counts
$n_{i'}$ do not decrease, and $n_i$ rises, because the owner $o_F\in(p,u)$ is an event of scale
$\lambda(d_F)\ge\beta_i$. So $\eta$ is violated only if \textup{(d1)} $x_0$ visits a level-$i$ line during $(p,u)$, or
\textup{(d2)} some $e_{i'}$ with $i'\ge i$ wraps in $(p,u]$.

\emph{Cost.} All expectations are over $\rho$ and the offsets; each violated edge is charged to one of the causes above,
and the classes may overlap, which only over-counts.
\begin{enumerate}[label=(\alph*),leftmargin=2em]
\item \emph{Edges at~$Z$.} A vertex has at most $2$ state edges, $2$ chords and $2k$ hops of~$H'$. Since
$\rho\bmod w$ is uniform, $\mathbb E\card Z=T/w$, and these edges cost at most $(4+2k)T/w\le\eps T/8$.
\item \emph{Descents.} By (M-a), they cost at most $T/K_0+\sum_i\#\{\text{events of scale}\ge\beta_i\}/K\le T/K_0+qNT/(2K)$,
because there are at most $\sum_v\sum_Fd_F\le NT/2$ events (\cref{lem:families}(e)).
\item \emph{(c1).} A wrap at time~$\tau$ of the level-$i$ counter of a strip~$S$ is charged only by chords at cells of~$S$
that span~$\tau$, at most $W_i-1$ of them by~(P1). At level~$0$ there are at most $T/K_0$ wraps in expectation,
costing at most $(w-1)T/K_0$. At level $i\ge1$ every entry follows a visit to a level-$i$ line, there are $T/W_i$ such
visits in expectation, and each entry is a wrap with probability at most $1/K$; so each level $i\ge1$ costs at most $T/K$.
\item \emph{(c2).} Let $e_{i'}$ wrap at~$\tau$. The chords charged to it have level $i<i'$ and span~$\tau$, so by
confinement $x_0(\tau)\in E(\gamma)\cup\{c\}\subseteq\sigma_{i+1}(c)\subseteq\sigma_{i'}(c)$: all of them lie at cells of the
level-$i'$ strip containing $x_0(\tau)$, and by~(P1) there are at most $W_{i'}-1$ of them. Events of scale at least
$\beta_{i'}$ are owners of families of height at least $2^{\beta_{i'}}$, so by \cref{lem:families}(e) there are at most
$NT/2^{\beta_{i'}+1}$ of them. The cost is at most
$\sum_{i'}W_{i'}NT/(2^{\beta_{i'}+1}K)\le\sum_{i'}(\mathrm{lo}_{i'}+w)NT/(2^{\beta_{i'}}K)\le
q(64qN^22^h/\eps+wN)T/K$.
\item \emph{(d1).} Charge $\eta$ to its first visit $\tau\in(p,u)$ to a level-$i$ line. The hops of window~$i$ that
span~$\tau$ lie in families of band at most $\beta_i+h$, so by the overlap bound of \cref{lem:families} there are fewer than
$N2^{\beta_i+h+1}$ of them. There are $T/W_i\le T/\mathrm{lo}_i$ visits to level-$i$ lines in expectation, so the cost is
at most $\sum_iN2^{\beta_i+h+1}T/\mathrm{lo}_i=\eps T/32$.
\item \emph{(d2).} An event~$o'$ of scale~$s$ raises the counts $n_{i'}$ with $\beta_{i'}\le s$, and a wrap of $e_{i'}$ at
$o'$ exposes the hops of windows at most~$i'$ that span~$o'$, whose bands are at most $\beta_{i'}+h$: fewer than
$N2^{\beta_{i'}+h+1}$ of them. Since the $\beta_{i'}$ are distinct integers, $\sum_{\beta_{i'}\le s}2^{\beta_{i'}}<2^{s+1}$,
so the cost is at most $\sum_{o'}N2^{h+2}2^{s(o')}/K\le N^22^{h+1}T/K$, using
$\sum_{o'}2^{s(o')}\le\sum_v\sum_Fd_F\le NT/2$.
\end{enumerate}
\emph{Total.} The edges at~$Z$ cost at most $\eps T/8$; the level-$0$ terms, $T/K_0+(w-1)T/K_0=wT/K_0$, at most
$\eps T/8$; all terms with denominator~$K$ together at most
$(qN+q+64q^2N^22^h/\eps+qwN+N^22^{h+1})T/K\le\eps T/8$; and (d1) at most $\eps T/32$. The sum is less than $\eps T/2$, and
some choice of $\rho$ and the offsets attains at most the expectation.
\end{informalproof}

The chords never see the digits of the bounded walks except through wraps, because the digit positions are static and the
event counts are global. A chord of a strip that $x_0$ has left is decided at the coarsest level whose line it crosses, by
the re-entry into its own strip there; above that level it is confined, which bounds its exposure in~(c2) by a strip
width rather than by the number of live chords. The hops of the bounded walks compare only head-$0$ digits at or above
their window, whose lines are far apart compared with their heights.

\begin{corollary}[Corollary M1: scale-ratio form]\label{cor:M1}
\notformalized
If all family heights of~$H'$ lie in $[d_{\min},d_{\max}]$, then $G$ has a cut of at most $\eps T$ edges and depth
\[
  O\bigl(k\,(N^2d_{\max}/d_{\min}+kN)^2/\eps^6\bigr)=\poly(k,\ell,1/\eps)\cdot(d_{\max}/d_{\min})^2 .
\]
In particular: (i) bounded walks whose families have height at most $d_{\max}$, together with any $x_0$, are cut at depth
$\poly(k,\ell,1/\eps)\,d_{\max}^2$; (ii) if all families except at most two per virtual head have their band in one window of
width~$h$, the depth is bounded in terms of $k$, $\ell$, $h$ and~$\eps$ alone; this covers zigzags of any amplitude,
even growing with~$T$.
\end{corollary}

\begin{informalproof}
Take $q=1$, $\beta_1=\lambda(d_{\min})$ and $h=\lambda(d_{\max})-\lambda(d_{\min})$, so that
$2^h\le2d_{\max}/d_{\min}$; then $K_0=O(k/\eps^2)$ and $K=O((N^22^h+kN)/\eps^2)$, and \cref{thm:M} gives depth less
than $K_0K^2$. For~(i), $d_{\min}\ge1$. For~(ii), the exceptional families are those truncated at the ends of $[0,T)$; with
them the bands are covered by at most $1+2N$ windows of width~$h$, and \cref{thm:M} applies with $q\le1+2N$. A zigzag
between two cells at distance~$A$ has owner-load one per side, and all its families have height~$A$ except the truncated
ones, so $h=0$.
\end{informalproof}

\paragraph{Why it matters.} Theorem~M allows one head to be anything: of unbounded owner-load, nested at every scale,
drifting, deep everywhere. So a depth-robust family must contain either two heads of unbounded owner-load, as the
deep-nesting pair does, or one such head together with bounded heads whose families occupy unboundedly many separated
scales (\cref{cor:where}). For machines this becomes \cref{cor:M-machine}.

\subsection{Where a counterexample to TCP must live}\label{sec:one-counter}

\begin{corollary}[Where depth-robustness lives]\label{cor:where}
\notformalized
Let $\mathcal W$ be a depth-robust family of $k$-tuples of walks. Then the owner-load of the set of all long hops is
unbounded on~$\mathcal W$. If moreover all walks but one have owner-load at most a fixed~$\ell$ on~$\mathcal W$, then for
every fixed~$h$ the least number of windows of width~$h$ covering the family bands of the bounded walks is unbounded
on~$\mathcal W$; in particular the ratio of their largest to their smallest family height is unbounded.
\end{corollary}

\begin{informalproof}
If the owner-load were bounded, \cref{thm:F} would give a single depth bound for each~$\eps$. If, for some~$h$, the bands
of the bounded walks of every member were covered by at most $q$ windows of width~$h$, \cref{thm:M}, applied with the
unbounded walk as~$x_0$, would do the same. A bounded ratio $d_{\max}/d_{\min}$ gives such a cover with $q=1$.
\end{informalproof}

The structure of walk DAGs that these theorems exploit is a stack structure. Call an \emph{$s$-stack DAG} a DAG on
$[0,N)$, with edges determined by their endpoints, consisting of the spine edges $t\to t-1$ and of \emph{arcs} $u\to p$ with
$u-p\ge2$, the arcs being partitioned into $s$ families; in each family no two arcs cross (there are no arcs $(p,u)$,
$(p',u')$ with $p<p'<u<u'$), and every vertex is the left endpoint of at most one arc of each family and the right endpoint
of at most one arc of each family. A cut is a set of edges, and its depth is the largest number of arcs on a path that avoids it. Let $\mathrm{SR}(s)$ be
the statement that for every $\eps>0$ there is $D$ such that every $s$-stack DAG on $[0,N)$ has a cut of at most $\eps N$
edges and depth at most~$D$.

\begin{proposition}[TCP is depth reduction for stack DAGs]\label{prop:stacks}
\notformalized
\begin{enumerate}[label=(\alph*)]
\item The walk DAG of $k$ walks, pauses allowed, is a $2k$-stack DAG.
\item Every $s$-stack DAG $H$ on $[0,N)$ is realized by $s$ walks of length $2N$ without pauses: every arc $(p,u)$ of
family~$i$ is the upper long hop $(2p,2u)$ of walk~$i$. Consequently a cut of size~$c$ and depth~$D$ of their walk DAG yields
a cut of~$H$ of size at most~$c$ and depth at most~$D$.
\end{enumerate}
Hence $\TCP(k)$ for walks without pauses implies $\mathrm{SR}(k)$, with the depth of $\TCP(k)$ at $\eps/2$, and
$\mathrm{SR}(2k)$ implies $\TCP(k)$ for walks with pauses; in particular the forms of TCP with and without pauses are
equivalent up to a factor~$2$ in the number of walks.
\end{proposition}

\begin{informalproof}
(a) The upper long hops of one walk do not cross, nor do its lower long hops (\cref{sec:prelim-walks}). Each time is the left
endpoint of at most one long hop of each walk and the right endpoint of at most one. A pair of times that is a long hop of
several walks is a single edge, kept in one of its families, and pauses are spine edges.

(b) For $0\le t\le N$ let $d_i(t)$ be the number of arcs of family~$i$ that contain~$t$ strictly inside, so
$d_i(0)=d_i(N)=0$. By the degree condition, $\abs{d_i(t+1)-d_i(t)}\le1$: $d_i(t+1)-d_i(t)$ is $1$ if $t$ is the left
endpoint of an arc and $0$ otherwise, minus $1$ if $t+1$ is the right endpoint of an arc and $0$ otherwise. The two
endpoints of an arc $(p,u)$ have the same value of~$d_i$: an arc $(p',u')$ of the family that contains~$p$ strictly inside
has $u'>u$, since $u'<u$ would cross and $u'=u$ would violate the degree condition, so it also contains~$u$, and
conversely. Every time strictly inside the arc has a larger value. Put $x_i(2t)=2d_i(t)$ and
$x_i(2t+1)=\min(x_i(2t),x_i(2t+2))+1$ for $t<N$. Consecutive even values differ by $0$ or~$\pm2$, so $x_i$ is a walk
without pauses. At $2p$ and $2u$ it is at $2d_i(p)$, and strictly between them it is at least $2d_i(p)+1$, so $(2p,2u)$
is an upper long hop of~$x_i$. A spine edge $t\to t-1$ maps to the state edges $2t\to2t-1\to2t-2$ and an arc $u\to p$ to the
hop edge $2u\to2p$; distinct edges of~$H$ have disjoint images. Given a cut~$F$ of the walk DAG, cut an edge of~$H$ if some
edge of its image lies in~$F$. This cut has at most $\card F$ edges, and a path of~$H$ that avoids it maps to a path that
avoids~$F$ with as many long hops as the path has arcs.

For the consequences, a $k$-stack DAG on $[0,N)$, padded with empty families, is realized by $k$ walks of length $2N$, and a
cut of $\eps N=(\eps/2)(2N)$ edges of their walk DAG pulls back; conversely~(a) applies to $k$ walks.
\end{informalproof}

In this language, Theorem~F says that $s$-stack DAGs are depth-reducible when the arcs open at any time in any family come
from boundedly many distinct owners, the owner of an arc being that of its hop in~(b). Forgetting the orientation of the
edges, $s$-stack DAGs are essentially the $s$-pushdown graphs of Galil, Kannan and Szemerédi~\cite{GKS1989separators}. Part~(a) of \cref{prop:stacks} is the
classical fact that a tape is two pushdown stores, so that the computation graphs of multitape machines, the objects of the
computation-graph method of Hopcroft, Paul and Valiant~\cite{HPV1977}, are $O(1)$-pushdown graphs; part~(b) is the converse,
a pushdown store being a tape whose head sits at the height of the stack. The stack view of tape computations was also used
by Paul, Pippenger, Szemerédi and Trotter~\cite{PPST1983} and by Santhanam~\cite{Santhanam2001}. $\mathrm{SR}(s)$ is the
depth analogue of the separator question of Galil, Kannan and Szemerédi, who showed that sublinear separators for
$3$-pushdown graphs, sublinear separators for $k$-page graphs, and subquadratic simulations of nondeterministic two-tape machines
by nondeterministic one-tape machines stand or fall together~\cite{GKS1989separators}. We found no work on reducing DAGs of bounded stack number, or walk DAGs
of bounded nesting, to bounded depth by deleting a linear fraction of the edges.

\begin{proposition}[Two stacks are planar]\label{prop:sr2}
\notformalized
$\mathrm{SR}(2)$ holds, with $D(\eps)=O(\eps^{-2})$.
\end{proposition}

\begin{informalproof}
Draw the vertices of a $2$-stack DAG in order on a line, the spine edges on the line, the arcs of one family above it and
those of the other below it. Arcs of one family do not cross, and arcs that share an endpoint can be drawn without
crossing, so the underlying graph is planar; its maximum degree is at most $2+2\cdot2=6$. By the planar separator theorem of
Lipton and Tarjan~\cite{LiptonTarjan1979}, applied recursively as in~\cite{LiptonTarjan1980}, some set of $O(N/\sqrt M)$
vertices leaves components of at most $M$ vertices. Cutting the at most six edges at each of these vertices cuts
$O(N/\sqrt M)$ edges, and a path that avoids the cut uses no vertex of the separator, except possibly as a path with no edge,
and so stays inside one component and has fewer than $M$ arcs. Taking $M$ of order $\eps^{-2}$ gives the claim.
\end{informalproof}

\paragraph{Why they matter.} \Cref{prop:sr2} recovers $\TCP(1)$, and \cref{prop:stacks} places the open part of TCP at
three or four stacks: $\mathrm{SR}(3)$ is the first open case, and if $\TCP(2)$ fails, as Conjecture~A$'$ with its walk-DAG
analogue predicts for the deep-nesting pairs, then $\mathrm{SR}(4)$ fails. From three stacks on, separators no longer help.
Galil, Kannan and Szemerédi constructed $3$-pushdown graphs whose separators have $\Omega(n/\log^{(s)}n)$ vertices, with
$\log^{(s)}$ the $s$-fold iterated logarithm, for every fixed~$s$~\cite{GKS1989pushdown}. Dvir and Wigderson obtained
$O(1)$-pushdown expanders from monotone expanders, answering a question of Galil, Kannan and
Szemerédi~\cite{DvirWigderson2010}, and Dujmović, Sidiropoulos and Wood built bounded-degree expanders with three-page book
embeddings from the constant-degree monotone expanders of Bourgain and Yehudayoff~\cite{DSW2016,BourgainYehudayoff2013};
each page of that embedding is a non-crossing matching, so with a spine added these expanders are $3$-stack DAGs.
Depth-robustness, however, needs arcs that span many scales, and within
boundedly many non-crossing families such arcs must nest; Theorem~F says that the nesting must then come from unboundedly
many distinct owners.

\begin{remark}[A candidate counterexample]\label{rem:umix-candidate}
Heuristic. Let $x_0=U(s_0,m,1)$, a nested hierarchy, and let $x_1$ be a sequential one: a walk made of V's, each a
segment of $2^{j+1}$ steps followed by a back stroke and a forward stroke of $2^j$ steps, with the exponents~$j$ running
through $s_m,\dots,s_0$ with frequencies proportional to $2^{-j}$, so that every scale gets the same time, as every level
of~$U$ does. Then $x_1$ has owner-load one per side and enters every cell at most three times. Theorem~F cuts $x_1$ alone
at depth $2^{O(1/\eps)}$ and \cref{lem:W} cuts $x_0$ alone at depth $O(\eps^{-2})$, while Theorem~M cuts the pair at depth
$O(4^m)\poly(1/\eps)$, with one window of width $h=m$. If the pair needs depth $2^{\Omega(m)}$ at some fixed~$\eps$, then
$\TCP(2)$ fails already for one deep head and one head of owner-load one. The fourth paper's forced-shape heuristics for the
deep-nesting pair apply with the small V's of~$x_1$ in the role of the second hierarchy; no lower-bound technique is
available, and packing certificates cannot certify depth-robustness~\cite[Proposition~5.14]{FourthPaper2026}.
\end{remark}
